% ============================================================
% Chapter 1: Introduction — Hilbert and His 23 Problems
% ============================================================

\chapter*{David Hilbert and His 23 Problems}
\addcontentsline{toc}{chapter}{Introduction: David Hilbert and His 23 Problems}

\section{The Man Behind the Problems}

David Hilbert (1862--1943) was born in K\"onigsberg, Prussia (now Kaliningrad, Russia), a city with a storied mathematical heritage---it was where Carl Jacobi had taught and where the young Hilbert would walk and talk mathematics with his friend Hermann Minkowski and their teacher Ferdinand von Lindemann.

Hilbert's early work was in invariant theory, where he proved a finite basis theorem that stunned the mathematical world. When Paul Gordan, the self-styled ``king of invariants,'' saw Hilbert's proof---a short, elegant argument where pages of computation had been the norm---he reportedly said: ``Das ist nicht Mathematik. Das ist Theologie.'' (``This is not mathematics. This is theology.'') But the proof was correct, and it opened a new way of thinking about existence proofs.

By 1900, Hilbert was already a prominent figure in G\"ottingen, one of the great mathematical centers of the world. He had moved from invariant theory to number theory, algebraic geometry, and the foundations of mathematics. He was a man of enormous vision, capable of seeing the deep connections between seemingly disparate areas of mathematics.

\section{The Second International Congress of Mathematicians}

The International Congress of Mathematicians (ICM) met in Paris in August 1900, timed to coincide with the opening of the Paris World's Fair. The Congress was a grand affair---a celebration of mathematics at the dawn of a new century. Henri Poincar\'e, the greatest mathematician of the age, delivered the opening address.

Hilbert was invited to give a lecture. His choice of topic was remarkable. Rather than presenting his own research, Hilbert decided to look forward. He would propose problems that should occupy mathematicians in the coming century.

\begin{quote}
\textit{``Who among us would not like to lift the veil behind which the future lies hidden and cast a glance at the coming advances of our science and the secrets of its development during future centuries? What new goals and what new methods will the new century reveal to us in the field of mathematics, whose old flora has been gathered so richly by the heroic exertions of previous centuries?''}
\par\hfill ---David Hilbert, 1900
\end{quote}

On the morning of August 8, 1900, Hilbert presented his list. It was not a complete list---he did not claim to have found all the important problems. ``Rather will it be the task of the future to follow the tracks he has blazed,'' he said of the problems of the past. ``Guided by the clear signs which they have left behind them, we may hope to gain a view of the whole of the yet unexplored domain of the mathematical sciences, and thus of the means whereby we may attack the problems which it contains.''

He then outlined twenty-three problems, spanning number theory, algebra, geometry, analysis, mathematical physics, and the foundations of mathematics. The problems were not chosen at random. They were carefully selected---each one was, in Hilbert's view, a doorway into a deep and important area of mathematics.

\section{The Problems: An Overview}

The twenty-three problems can be grouped into several broad areas:

\begin{enumerate}
    \item \textbf{Number Theory and Algebra} (Problems 1, 8, 9, 10, 11, 12, 13, 14): Questions about prime numbers, Diophantine equations, algebraic number fields, and the solvability of equations by radicals.

    \item \textbf{Geometry} (Problems 3, 4, 5, 6, 16, 18): Questions about volume, the concept of distance, continuous groups, the axioms of geometry, and tessellations of space.

    \item \textbf{Analysis} (Problems 2, 15, 17, 19, 20, 21, 22, 23): Questions about the continuum, calculus of variations, boundary value problems, and the monodromy theorem.

    \item \textbf{Mathematical Physics} (Problems 6, 19): The relationship between physics and mathematics, and the problem of regularity of solutions to equations of mathematical physics.

    \item \textbf{Logic and Foundations} (Problems 2, 7): The consistency of arithmetic and the nature of irrational numbers.
\end{enumerate}

The following table gives a bird's-eye view of all 23 problems and their current status:

\begin{center}
\small
\renewcommand{\arraystretch}{1.2}
\begin{tabular}{clll}
\toprule
\textbf{\#} & \textbf{Problem (short title)} & \textbf{Year Solved} & \textbf{Status} \\
\midrule
1 & Continuum Hypothesis & 1963 & Independent \\
2 & Consistency of Arithmetic & 1931 & Resolved \\
3 & Equality of Volumes (Tetrahedra) & 1901 & Solved \\
4 & Shortest Path = Straight Line & 1900s & Solved \\
5 & Continuous Groups = Manifolds & 1920s & Solved \\
6 & Axioms of Physics & --- & Partially Solved \\
7 & Irrationality of $e^\pi$ & 1934 & Solved \\
8 & Goldbach Conjecture & --- & Unsolved \\
9 & Reciprocity Laws & 1920s & Solved \\
10 & Decision for Diophantine Equations & 1970 & Solved \\
11 & Quadratic Forms & 1930s & Partially Solved \\
12 & Kronecker's Theorem & 1920s & Solved \\
13 & Solving 7th Degree by Radicals & 1900s & Solved \\
14 & Constructive Algebraic Invariants & 1950s & Partially Solved \\
15 & Schubert's Enumerative Calculus & --- & Partially Solved \\
16 & Topology of Algebraic Curves & --- & Partially Solved \\
17 & Representation by Quadratic Forms & 1920s & Solved \\
18 & Tessellation of Space & 1910--11 & Solved \\
19 & Regularity of PDE Solutions & 1904--57 & Partially Solved \\
20 & Boundary Value Problems & --- & Partially Solved \\
21 & Monodromy Problem & 1950s & Solved \\
22 & Uniformization & 1907--13 & Solved \\
23 & Calculus of Variations & 1900s+ & Partially Solved \\
\bottomrule
\end{tabular}
\end{center}

\section{Why These Problems Matter}

Hilbert's problems are not just historical artifacts. They remain central to modern mathematics. Several of the deepest results of 20th-century mathematics---G\"odel's incompleteness theorems, the proof of the Weil conjectures, the resolution of the Poincar\'e conjecture (not one of Hilbert's, but related in spirit), the classification of finite simple groups---were motivated directly or indirectly by Hilbert's vision.

More importantly, the problems illustrate what mathematics \emph{is}: a living, evolving discipline where the deepest questions of one generation become the foundational theorems of the next.

\section{A Guide for the Reader}

This book proceeds as follows. Each of the next 23 chapters addresses one of Hilbert's problems. The chapters are written to be self-contained---you may read them in any order. Each chapter includes:

\begin{itemize}
    \item \textbf{Historical Context}: What was known in 1900 and why the problem mattered.
    \item \textbf{The Problem}: Stated in modern mathematical language.
    \item \textbf{Key Developments}: The major ideas and results that addressed the problem.
    \item \textbf{Current Status}: Where the problem stands today.
    \item \textbf{Exercises}: Problems to test your understanding.
\end{itemize}

No single book can teach all the mathematics needed to fully appreciate every one of Hilbert's problems. But we can convey the essential ideas, and we can inspire the reader to seek out more.

\section*{Exercises}

\begin{enumerate}
    \item Research and write a one-page biography of David Hilbert. What were his major contributions to mathematics beyond the 23 problems?

    \item Hilbert said, ``Wir m\"ussen wissen. Wir werden wissen.'' (``We must know. We will know.'') What did he mean, and how does this relate to the problems he proposed?

    \item Choose three of the 23 problems that you find interesting after reading the table above. For each, write one sentence explaining why you chose it.
\end{enumerate}
